\section{System Design}

\subsection{Backend-Centered Pipeline}
The backend AI pipeline represents the core of Examiner Coach. The pipeline receives spoken
feedback, obtains a transcript, retrieves relevant educational evidence,
evaluates the transcript against a rubric, and supports follow-up coaching.

The pipeline can be summarized as follows:
\begin{enumerate}
    \item the examiner records spoken feedback for an OSCE station;
    \item the audio is transcribed by the configured speech model;
    \item the transcript is normalized where needed for retrieval;
    \item criterion-aware retrieval queries are constructed;
    \item educational feedback literature is retrieved from the vector store;
    \item the LLM evaluates the transcript using the rubric and retrieved
    evidence;
    \item the structured result is validated and resolved into the requested
    display language;
    \item optional follow-up coaching chat reuses the transcript, evaluation, and
    retrieved evidence.
\end{enumerate}

The frontend is therefore intentionally not the scientific focus of this
report. It acts as an interaction layer for recording and displaying results,
whereas the research contribution is the modelling and orchestration of the
backend evaluation and coaching process.

The design is shaped by the problem definition: the system must be available for
individual practice, return feedback immediately enough to support rehearsal,
and evaluate the examiner's response against explicit criteria rather than
providing only generic comments. It must also keep the data flow inspectable so
that privacy implications can be discussed and controlled in an institutional
medical education context. In this sense, the project can be understood as a
small applied step toward the KISSKI vision: providing AI methods for sensitive
health-related infrastructures in a way that is usable for real institutional
workflows rather than only for isolated demonstrations.

\subsection{Knowledge Base and Retrieval}
Educational feedback material is ingested into a local ChromaDB vector store.
Documents are chunked, embedded using a multilingual embedding model respecting the bilingual nature of both the knowledge base and the feedback transcripts, and
retrieved at runtime based on the feedback transcript and evaluation criteria.
Retrieval can be performed directly or through HyDE-based hypothetical-passage
retrieval.

The knowledge base gives the system a way to ground feedback evaluation in
educational guidance instead of relying only on the LLM's internal training
data. This is especially important because the task is not generic text
classification. The system must distinguish feedback that is merely positive
from feedback that is specific, formative, and behaviorally actionable. The
retrieved context provides concepts and examples that help the evaluation prompt
remain aligned with feedback education principles.

\begin{table}[ht]
    \centering
    \small
    \begin{tabular}{L{0.20\linewidth}L{0.32\linewidth}L{0.36\linewidth}}
        \toprule
        Module & Implementation in Examiner Coach & Purpose \\
        \midrule
        Indexing & Document ingestion, chunking, embeddings, ChromaDB storage
        & Prepare educational evidence for later retrieval. \\
        Pre-retrieval & Transcript normalization and criterion-aware query
        construction & Improve retrieval relevance across German and English
        feedback. \\
        Retrieval & Direct vector search or HyDE-based retrieval & Select
        candidate evidence chunks from the knowledge base. \\
        Post-retrieval & Filtering, reranking, and final evidence formatting
        & Prefer practical feedback guidance over generic or weakly related
        chunks. \\
        Generation & Structured LLM prompt using transcript, evidence, and
        rubric & Produce criterion scores, suggestions, quotes, and summary. \\
        Validation & JSON parsing and schema-based result assembly & Keep
        output machine-readable and suitable for coaching reuse. \\
        \bottomrule
    \end{tabular}
    \caption{Mapping of Examiner Coach to a modular RAG architecture.}
    \label{tab:modular-rag-mapping}
\end{table}

\subsection{Evaluation and Coaching Flow}
The evaluation flow combines the transcript, retrieved evidence, and a
structured rubric prompt. The LLM returns criterion-level scores, suggestions,
and quotes where applicable. The backend then parses the response, computes the
overall score, validates the structured result, and resolves it into the
requested display language. The coaching flow reuses the transcript,
evaluation result, and retrieved evidence to answer follow-up questions.

The coaching component extends the evaluation from a static score report into a
learning interaction. After receiving an evaluation, the examiner can ask why a
criterion was scored in a particular way or request alternative wording. This is
important pedagogically because the score alone does not teach the examiner how
to improve. The coaching response can refer back to the transcript, the scored
criteria, and the evidence base, thereby connecting system output to concrete
feedback practice.
